\newpage % Rozdziały zaczynamy od nowej strony.
\section{Praefatio}
\lipsum[1] \cite{goossens93}
\begin{figure}[!h]
    \label{fig:tradycyjne-logo-pw}
    \centering \includegraphics[width=0.5\linewidth]{logopw.png}
    \caption{Tradycyjne godło Politechniki Warszawskiej}
\end{figure}
\lipsum[2-3]
\begin{figure}[!h]
	\label{fig:nowe-logo-pw}
	\centering \includegraphics[width=0.5\linewidth]{logopw2.png}
	\caption{Współczesne logo Politechniki Warszawskiej}
\end{figure}
\lipsum[4-6]
Każdy post posiada zbiór cech, który może przyjmować jedną z dwóch możliwych wartości – 1, gdy dany użytkownik udostępnił dany post lub 0, gdy dany użytkownik nie udostępnił danego postu.

$$\forall\ i\in I\ ,\ u\in U\ \ \ \exists\ x_{iu}=f\left(i,u\right)$$

$$f\left(i,u\right)\ =\ 1,gdy użytkownik u udostępnił post i  \\       0,gdy użytkownik u nie udostępnił postu i $$



że istnieje jeszcze duża potrzeba i możliwości w temacie badania i wykrywania nierzetelnych informacji. Informacje rozprzestrzeniają się bardzo szybko dzięki użyciu mediów społecznościowych. Przedstawiona w tej pracy metoda, mimo stosunkowo wysokich wyników, nie jest jednak możliwa do użycia w realnym środowisku. Pierwszą przeszkodą jaką należało by pokonać jest dostęp do danych który nie był by tak ograniczony. 

nie wyczerpuje tematu ale daje doskonaly fundament do kolejnych prac.. udalo sie udowodnic ze to jest obiecujaca sciezka aby dalej ta sciezke formatowac. 

przeczytac tytuly czy czy nie ma literowek.
  techniczny opis analizy danych 

rozmowa z recenzentem na teams bedzie wazna.  Dominik Ruszko. poprosic o rozmowe online. 

